\section{The theory, by field type}\label{sec:theory}

% Survey section (draft 2026-10-15). Sources: LITERATURE.md; every
% statement below is at abstract level until the full papers are read.

\subsection{What a decomposition costs}

Every index-calculus algorithm for elliptic curves in the literature
spends its time in three places: finding about $\abs\Fb$ relations (each
a decomposition of a random combination $aG+bQ$ over $\Fb$), solving the
sparse linear system over $\ZZ/r$, and, in some variants, a descent of
the target. The second and third are well understood. The first is the
crux: a uniformly random point decomposes over a base of size $N$ into $m$
summands with probability about $N^m/(m!\,r)$, so relations need about
$m!\,r/N^{m-1}$ trials, and each trial costs whatever the decomposition
oracle costs. Every result below is a statement about one of three things:
how large $N$ can be while membership stays algebraically testable, how
cheap the oracle is, and how the two trade against the linear algebra.

\subsection{Extension fields $\FF_{q^k}$ with $k>1$ fixed or small}

Gaudry~\cite{gaudry2009} takes $\Fb$ to be the points with abscissa in the
base field $\FF_q$ and decomposes into $m=k$ summands by Weil restriction of
$\Sem{k+1}$, a system of $k$ equations in $k$ unknowns over $\FF_q$. For
fixed $k$ this gives $\tilde O(q^{2-2/k})$, asymptotically below rho's
$q^{k/2}$ for $k\ge3$; the constant hidden in the Gr\"obner solve grows
very quickly with $k$, which is what makes the crossover large in
practice (to be quantified from the paper).
Diem~\cite{diem2011} extends the analysis to families where $k$ grows,
proving subexponential complexity over sequences of fields with suitably
balanced $q$ and $k$. Joux and Vitse's work on decomposition and cover
attacks over small extension fields is to be added here once read. FGHR~\cite{fghr2014} exploit the action of $S_m$ and of
the 2-torsion to replace $\Sem{m+1}$ by its invariant form, cutting the
Gr\"obner work by large factors; our measurement at $k=4$
(\S\ref{sec:gaudry}) is a direct instance.

\subsection{Binary fields $\FF_{2^n}$ with $n$ prime}

Prime $n$ leaves no subfield, so the factor base is taken to be the points
whose abscissa lies in an $\FF_2$-vector subspace $V$ of dimension $\ell$.
Weil restriction then gives a quadratic-to-low-degree Boolean system whose
difficulty is governed by its degree of regularity. The whole debate about
binary curves is a debate about that degree.

\paragraph{The subexponential claims.} Petit and
Quisquater~\cite{petit2012} conjectured, from experiments and heuristic
evidence, that these systems have degrees of regularity only slightly above
the degrees of their equations (the \emph{first fall degree assumption}),
and deduced a heuristic $2^{c\,n^{2/3}\log n}$, $c<2$, for a variant of
Diem's algorithm, beating generic methods only for $n\gtrsim2000$.
Semaev~\cite{semaev2015} split summation polynomials into cubic Boolean
systems, argued for regularity degree at most 4 under the same kind of
assumption, and obtained a heuristic $2^{c\sqrt{n\ln n}}$ with
$c\approx1.69$; its first version concluded that several FIPS binary curves
fall faster than Pollard rho, a conclusion its revision says does not hold
under the weaker assumption that the regularity degree is
$o(\sqrt{n/\ln n})$.

\paragraph{The evidence against.} Shantz and Teske~\cite{shantz2013}
tested the Petit--Quisquater heuristics experimentally. Huang, Kosters and
Yeo~\cite{hky2015} introduced the \emph{last fall degree}, independent of
the monomial order, gave computational evidence against the first fall
degree assumption in characteristic~2, and constructed summation-polynomial
Weil descent systems where it is unlikely to hold. Galbraith and
Gebregiyorgis~\cite{galbraith2014} improved the variable choice (binary
Edwards curves), found symmetry-breaking factor bases and SAT solvers
inconclusive, and concluded that rho remains much faster for prime-degree
binary fields of reasonable size. The Galbraith--Gaudry
survey~\cite{galbraith2016} reaches the same verdict for characteristic-two
curves over prime-degree fields.

\paragraph{Where it stands.} No proof or refutation of a subexponential
algorithm for prime-degree binary curves exists. The open quantity is the
solving degree of Weil-descended $\Sem{m+1}$ as $n$ and $m$ grow; recent
work supplies diagnostics (saturation-exponent bounds, multi-graded
regularity, tests for where systems stop looking generic) but no bound for
these systems. Our own registered refutation-degree law at $m=3$
(\S\ref{sec:groebner}) is one data point and ended inconclusive.

\subsection{Subfield (Koblitz) curves}

When the curve is defined over $\FF_2$ the Frobenius acts on the group,
and Galbraith, Granger, Merz and Petit~\cite{ggmp2020} show that a
Frobenius-invariant factor base saves a factor $\approx n$ in relations and
$\approx n^2$ in linear algebra---more than rho's $\sqrt{2n}$---while
concluding that Koblitz-curve security is not threatened at deployed sizes.
Our Koblitz work (\S\ref{sec:koblitz}) builds on this.

\subsection{Prime fields}

Over $\FF_p$ there is no subfield and no subspace, and no factor base is
known for which membership is both dense enough and algebraically testable.
Semaev's original small-abscissa base leads to systems of very high degree.
Petit, Kosters and Messeng~\cite{pkm2016} construct bases from rational maps
(including the Kummer tower $\mu_{2^t}$) that give the summation system
extra structure; Amadori, Pintore and Sala~\cite{aps2018} build the base
from the instance and need only one relation search, improving on both;
McGuire and Mueller~\cite{mcguire2017} avoid Gr\"obner bases altogether by
fast evaluation of $\Sem m$ through iterated resultants, and in a second
variant by summing $m-1$ factor-base points and testing membership of the
result---the same mechanism as the two-summand root-finding oracle of
\S\ref{sec:rootfinding}. All of these remain exponential and slower than
rho; a 2024 analysis again concludes that no subexponential prime-field
method is known.

\subsection{Lower bounds}

Generic-group lower bounds bound everything that does not use curve
structure. With preprocessing, Corrigan-Gibbs and
Kogan~\cite{cgk2018} show $S\,T^2=\tilde\Omega(\varepsilon N)$ for $S$ bits
of advice and online time $T$; this is the bound against which the
retained-set-up (online) results of \S\ref{sec:koblitz} must be read. A
2026 structured-generic-group bound shows that an algorithm exploiting
structure on only a $\delta$ fraction of elements needs
$\Omega(\min(\sqrt q,1/\delta))$ time---for index calculus, the
decomposable set must be dense \emph{and} cheaply testable.
